\section{Modelo de Dominio}

El dominio del problema se organiza alrededor del concepto de transacción de pago como entidad central, la cual es procesada por una pasarela de pago a través de métodos de pago específicos y genera notificaciones asíncronas hacia la aplicación del comercio. A continuación se describen los conceptos principales del dominio y sus relaciones.

\subsection{Entidades del dominio}

\renewcommand{\arraystretch}{1.2}
\begin{longtable}{|l|p{0.76\textwidth}|}
  \hline
  \textbf{Entidad} & \textbf{Descripción} \\
  \hline
  \endfirsthead
  \hline
  \textbf{Entidad} & \textbf{Descripción} \\
  \hline
  \endhead
  \hline
  \endfoot
  \caption{Entidades del modelo de dominio de Kit Pagos Colombia.}
  \label{tab:entidades-dominio}\\
  \endlastfoot
  Transacción & Unidad atómica de intercambio económico entre un comprador y un comercio, procesada a través de una pasarela de pago. Tiene un identificador único, un monto, una moneda, un estado, una marca temporal de creación, el estado nativo de la pasarela y, cuando existe, un motivo de rechazo. \\
  \hline
  Estado de transacción & Representa la condición actual de una transacción en el ciclo de vida del pago. Los estados normalizados del framework son: \texttt{APPROVED}, \texttt{DECLINED}, \texttt{PENDING}, \texttt{EXPIRED}, \texttt{VOIDED} y \texttt{ERROR}. \\
  \hline
  Resultado de pago & Unión que distingue dos desenlaces de la creación de un pago: una transacción resuelta (\texttt{TRANSACTION}) o una redirección requerida (\texttt{REDIRECT\_REQUIRED}) con la URL a la que hay que enviar al pagador, el identificador de la pasarela y el estado nativo. \\
  \hline
  Pasarela de pago & Proveedor externo que procesa las transacciones de pago. Cada pasarela define sus propias convenciones de API, autenticación, formatos de respuesta y mecanismos de notificación. Las pasarelas soportadas son Wompi, Rapyd (PayU GPO), Mercado Pago y Kushki. \\
  \hline
  Adaptador & Componente del framework que traduce el contrato normalizado del puerto hacia las convenciones específicas de una pasarela determinada. Existe un adaptador por cada pasarela soportada. \\
  \hline
  Puerto & Interfaz abstracta del framework que define el contrato que cada adaptador debe cumplir. El puerto especifica las operaciones disponibles (crear pago, consultar estado, listar bancos PSE y validar webhook) sin conocimiento de ningún proveedor específico. \\
  \hline
  Método de pago & Mecanismo a través del cual el comprador efectúa el pago. El framework soporta tarjeta de crédito o débito, a partir de un token emitido por la pasarela, y transferencia bancaria PSE. \\
  \hline
  Banco PSE & Entidad financiera habilitada para cobrar por PSE. Se identifica con el código que exige cada pasarela y, cuando está en el catálogo, con el código de compensación de ACH Colombia. \\
  \hline
  Webhook & Notificación asíncrona que una pasarela envía hacia una URL registrada por el comercio para informar sobre cambios en el estado de una transacción. Cada pasarela define su propio formato y mecanismo de firma. \\
  \hline
  Evento de webhook & Representación normalizada de una notificación válida. Incluye el tipo de evento, la pasarela, el identificador de la transacción, el nuevo estado y, cuando la firma lo cubre, la referencia de la orden. \\
  \hline
  Motivo de rechazo & Explicación de un pago \texttt{DECLINED}. Incluye el código nativo de la pasarela y una categoría normalizada: \texttt{INSUFFICIENT\_FUNDS}, \texttt{INVALID\_CARD\_DATA}, \texttt{CARD\_RESTRICTED}, \texttt{FRAUD\_SUSPICION}, \texttt{NETWORK\_ERROR}, \texttt{AUTHENTICATION\_FAILED}, \texttt{DUPLICATED\_PAYMENT} o \texttt{UNKNOWN}. \\
  \hline
  Error normalizado & Representación unificada de una condición de fallo durante el procesamiento de una operación, independiente del formato de error específico de cada pasarela. Incluye un código de error, un mensaje descriptivo y la pasarela que originó el fallo. \\
  \hline
  Comercio & Entidad que ofrece bienes o servicios e integra el framework para aceptar pagos digitales. Posee credenciales de autenticación para cada pasarela que desea utilizar. \\
  \hline
  Desarrollador & Persona técnica que consume el framework para integrar capacidades de pago en la aplicación del comercio. Es el usuario principal del SDK, del módulo de navegador y de la API de Simulación. \\
  \hline
\end{longtable}

\subsection{Relaciones entre entidades}

\begin{itemize}
  \item Un Comercio se autentica ante una o varias Pasarelas de pago mediante credenciales específicas por proveedor.
  \item Una Transacción es creada por un Comercio y procesada por exactamente una Pasarela de pago a través de un Método de pago.
  \item Una Transacción tiene exactamente un Estado de transacción en cada momento de su ciclo de vida.
  \item Una Transacción \texttt{DECLINED} puede tener un Motivo de rechazo. Un estado \texttt{ERROR} representa una falla técnica y no un rechazo del banco.
  \item Un pago con PSE referencia un Banco PSE cuyo código es válido en la pasarela que procesa la transacción.
  \item Una Pasarela de pago envía cero o más Webhooks hacia el Comercio a lo largo del ciclo de vida de una Transacción; el SDK los valida y los convierte en un Evento de webhook.
  \item El framework expone un Puerto que es implementado por uno o más Adaptadores, uno por cada Pasarela soportada.
  \item La creación de un pago produce un Resultado de pago: una Transacción resuelta o una redirección pendiente.
  \item Cada operación fallida genera un Error normalizado que abstrae el error original de la Pasarela.
\end{itemize}
